\chapter{Networking and Session Semantics}
\label{ch:networking}

CNA's \cnans{Microsoft::Xna::Framework::Net} namespace --- five enums and eighteen classes,
mirroring the shape this book's ecosystem chapter already summarized --- is one of the
clearest illustrations in this book of the gap between ``the API surface exists'' and ``the
API surface does something real.'' Real networking is backed by ENet (vendored directly, as
already noted earlier in this book), wrapped by an internal \texttt{ENetBackend}.

\section{NetworkSession: caller-owned, and genuinely capable in ways FNA's stub is not}

\cnaclass{NetworkSession} is explicit about its own ownership model: \texttt{Create()} /
\texttt{Find()} / \texttt{Join()} / \texttt{JoinInvited()} all return a caller-owned
\cnaclass{NetworkSession}\texttt{*}; \texttt{Dispose()} releases the transport and every gamer the
session owns, but does not delete the session object itself --- the caller must do that
separately, matching the conventional \texttt{IDisposable} split between resource release and
memory ownership. Two properties are worth calling out specifically because they are real in
CNA in a way they are not in FNA at all: \texttt{AllowHostMigration}, where every surviving
peer independently computes the identical new host (the lowest remaining wire ID, requiring
no election round-trip) and reconnects via real LAN discovery; and
\texttt{SimulatedLatency} / \texttt{SimulatedPacketLoss}, which really do delay or
probabilistically drop \texttt{AppData} traffic on the receiving side, scoped specifically to
application data rather than session-management traffic or host-relay hops. In FNA, both of
these properties exist only as stored, inert values with no behavior behind them at all ---
here, they are genuinely functional network-testing tools.

A subtle but load-bearing detail belongs right alongside those two: real XNA's own
\texttt{GamerJoined} is documented to replay itself immediately upon subscription for every
gamer already in the session at that point.  CNA now reproduces that contract through
\texttt{EventHandler<T>::SetReplayHook}: adding a handler synchronously invokes it once for
every gamer already in \texttt{AllGamers}, including the initial local gamer(s) established by
\texttt{Create()} / \texttt{Join()} / \texttt{JoinInvited()}.  Mid-session joins still enter
the ordinary network-event queue and therefore require the normal per-frame \texttt{Update()}
pump; an extra \texttt{Update()} immediately after subscribing is neither required nor a way
to obtain a second copy of the initial events.

\begin{lstlisting}
// Worked example: hosting a session and correctly receiving the local
// gamer's initial GamerJoined event through replay-on-subscribe.
NetworkSession* session = NetworkSession::Create(
    NetworkSessionType::SystemLink, /* maxLocalGamers */ 1, /* maxGamers */ 8);
    // the 8 above is NOT actually honored -- see the worked example below

session->GamerJoined.Add(
    [](System::Object* sender, const GamerJoinedEventArgs& args)
    {
        // This runs synchronously during Add() for gamers already present;
        // later remote joins arrive from the ordinary Update() pump.
    });

// ... in the game loop, every frame:
session->Update();
\end{lstlisting}

\section{Host migration, worked through in full}

\texttt{AllowHostMigration}'s ``lowest remaining wire ID, no election round-trip'' claim above
is worth unpacking into the actual algorithm, since it is a genuinely interesting piece of
distributed-systems design for a library this size. The detection point is the same
disconnect handler that already runs when any peer drops; when the disconnecting peer is
specifically the host, \texttt{AllowHostMigration == false} preserves today's exact behavior
(the session ends immediately), while \texttt{true} branches into a dedicated migration
helper instead. Every surviving peer already knows the full roster --- gamertag, wire ID, and
\texttt{IsHost} --- from the handshake messages exchanged when it originally joined, so no new
negotiation message is needed at all: each survivor independently computes
\texttt{new\_host = min(every other known wire ID)}, and if that ID belongs to one of its own
local gamers, it promotes itself; otherwise, it knows it must reconnect to whichever peer now
owns that ID. Both outcomes share the same cleanup step first --- every currently-known remote
\cnaclass{NetworkGamer} is removed via a real \texttt{GamerLeave} event, and all of the
session's internal wire-ID bookkeeping is cleared and reset to zero, so the next real handshake
cannot collide with stale entries left over from the dead host's own numbering.

\begin{lstlisting}
// Worked example: subscribing to the real HostChanged event that fires
// once a promoted or reconnected peer's migration actually completes.
session->HostChanged.Add(
    [](System::Object* sender, const HostChangedEventArgs& args)
    {
        // args.getOldHostProperty() / getNewHostProperty() -- both real
        // NetworkGamer pointers; getNewHostProperty() is only raised once
        // a live NetworkGamer for the new host genuinely exists again.
    });
\end{lstlisting}

The reconnecting side (every survivor that is \emph{not} promoted) has no direct channel to
any other peer in this star topology, and the old host is gone, so the only way to find the
new host's address is the same real LAN discovery \cnaclass{NetworkSession::Find()} already
uses --- matched against the expected new host's already-cached gamertag --- followed by an
ordinary reconnect through the exact same \texttt{ClientHello} / \texttt{ServerWelcome} handshake
a fresh \texttt{Join()} performs. This is a deliberate, explicitly confirmed scope choice worth
stating precisely: it is a full reconnect, not a live migration of existing sockets, so a
reconnecting peer's remote-gamer roster is genuinely rebuilt from scratch (fresh
\cnaclass{NetworkGamer} identities, real \texttt{GamerLeave} followed by real \texttt{GamerJoin}
events) rather than preserved across the promotion.  Discovery is retried for at most three
150~ms search windows to tolerate the promoted peer registering at nearly the same instant or
an OS delivering a shared-port discovery datagram to another socket. If none finds the cached
gamertag, migration falls through to the same immediate-end path
\texttt{AllowHostMigration == false} always takes.

\section{NetworkGamer: real per-machine identity where FNA hardcodes it}

\cnaclass{NetworkGamer} / \cnaclass{LocalNetworkGamer}'s \texttt{Id} and \texttt{IsHost} are
real, cross-machine-consistent values in CNA --- a genuine improvement over FNA's own
documented upstream limitation, where \emph{every} gamer's \texttt{Id} is hardcoded to zero
and \emph{every} gamer's \texttt{IsHost} is hardcoded to true, regardless of which machine is
asking. \texttt{RoundtripTime} is wired to the real underlying ENet peer's own native
round-trip-time tracking, not estimated or faked.

\section{Real versus stub, by NetworkSessionType}

This is the single most useful table in this chapter, because it tells you directly whether
a given multiplayer feature in a game you are porting (this book's migration chapter returns
to this) will actually work:

\begin{center}
\begin{tabular}{ll}
\toprule
\textbf{\cnaclass{NetworkSessionType}} & \textbf{Status} \\
\midrule
\texttt{SystemLink} & Real: hosting, joining, LAN discovery, AppData relay, disconnect \\
                     & handling, host migration --- all genuinely networked over ENet \\
\texttt{Local} & Synthetic lifecycle only: \texttt{SendData} delivers nothing \\
\texttt{LocalWithLeaderboards} & Same non-delivery, plus locally persisted \\
                                & leaderboards/achievements from Chapter~\ref{ch:gamerservices} \\
\texttt{PlayerMatch} / \texttt{Ranked} & Documented stub --- no matchmaking backend exists \\
                                        & to connect to, not an artificially withheld feature \\
\bottomrule
\end{tabular}
\end{center}

Session invites are stubbed for the identical reason \texttt{PlayerMatch} / \texttt{Ranked}
are: there is no online service behind them to implement against, not a deliberately
Xbox-Live-exclusive feature CNA chose not to attempt.

For all four synthetic types, no port is bound, \texttt{ConnectToHost} is a no-op, and
\texttt{SendData} leaves no local data available.  Discovery differs by type:
\texttt{Find(Local, ...)} rejects \texttt{Local} with \texttt{ArgumentException}, while
\texttt{LocalWithLeaderboards}, \texttt{PlayerMatch}, and \texttt{Ranked} return an empty
collection without the 150~ms network wait. Their useful behavior is session lifecycle:
\texttt{StartGame}, \texttt{EndGame}, state changes, and events. Only \texttt{SystemLink}
passes \texttt{RealNetworkingEnabled}.

\subsection{What remains a parity stub even in a real \texttt{SystemLink} session}

The functional ENet path should not be read as blanket completion of every XNA networking
member.  Voice transport is absent: \texttt{LocalNetworkGamer::EnableSendVoice()} and
\texttt{SendPartyInvites()} are no-ops, while \texttt{HasVoice}, \texttt{IsTalking},
\texttt{IsMutedByLocalUser}, \texttt{IsGuest}, and \texttt{IsPrivateSlot} remain their default
false values.  \texttt{NetworkMachine::RemoveFromSession()} always throws
\texttt{NotImplementedException}.  The three arbitrated/unarbitrated/TrueSkill leaderboard
events and the static \texttt{InviteAccepted} event are declared but never raised.  Invitations
have no transport or service behind them; \texttt{JoinInvited()} merely constructs the same
synthetic \texttt{PlayerMatch} lifecycle object described above.  These limits are independent
of the real discovery, handshake, application-data, state-broadcast, and host-migration paths.

\section{Real LAN discovery, concretely}

\cnaclass{AvailableNetworkSession} / \texttt{AvailableNetworkSessionCollection} are populated by
genuine ENet LAN discovery replies over a well-known UDP port, with broadcast and loopback
fallback and OS-assigned ephemeral ports for the actual game-session transport once a session
is joined. Several accessor methods on \cnaclass{AvailableNetworkSession} --- resolving the
real connect address, port, and session type --- exist purely because CNA's own
\texttt{Find()} genuinely discovers sessions over the network; they have no FNA counterpart
at all, since FNA's own \texttt{Find()} never discovers anything real to begin with.
\cnaclass{QualityOfService} correspondingly has two construction paths: one matching FNA's
own acknowledged-incomplete, no-measurement stub exactly (byte-for-byte, for API parity), and
a second that computes a real round-trip-time measurement from the discovery query/reply
exchange itself --- bandwidth remains unmeasured in both cases, since no established peer
connection exists yet at discovery time to measure it against.

The discovery protocol uses UDP port 61190, version~1, and a 150~ms search window. A finder
sends both broadcast and an explicit loopback query, then deduplicates replies by connect port.
Ordinary hosts use an OS-assigned ENet port; Emscripten's fixed host port is 61191, but raw UDP
discovery is a permanent browser no-op. Session transport allows 31 peers, two ENet channels,
and 64 queued pre-handshake application sends.

\subsection{Wire opcodes, topology, and hostile input}

The complete application protocol is small: \texttt{ClientHello=0x01},
\texttt{ServerWelcome=0x02}, join/leave broadcasts \texttt{0x03}/\texttt{0x04},
\texttt{0x05} reserved for an unimplemented host-change broadcast,
\texttt{StateChangeBroadcast=0x06}, and \texttt{AppData=0x10}. The topology is a star: clients
talk to the host and the host relays application data to its target. Host migration instead
promotes locally or rediscovers the new host by gamertag, retrying three 150~ms searches.

Untrusted discovery packets are version-checked before deeper parsing. Property indices must
be non-negative and below 256, preventing a crafted roughly two-billion-entry growth loop.
List counts are one byte and encoding throws above 255 rather than wrapping. The threat model
includes a custom ENet client speaking this inferable wire format; a non-host peer sending a
host-only broadcast is logged and disconnected, as is a completed peer that repeats its hello.

\section{Real bugs found and fixed}

Several genuine, consequential bugs were found and fixed in this namespace's own history,
worth naming because of what each reveals. \cnaclass{NetworkSession}'s own
\texttt{Dispose()} could be called twice and use-after-free as a result --- fixed to be
idempotent. The async \texttt{Begin*} family's callbacks were stored but, for a period, never
actually invoked at all --- fixed to fire exactly once, correctly handling the case where the
callback itself re-entrantly triggers another async call. Every \cnaclass{NetworkGamer} /
\cnaclass{LocalNetworkGamer} a session ever created used to leak permanently, because nothing
in the session's own lifecycle ever deleted them --- fixed by having \cnaclass{NetworkSession}
take real ownership and free them in bulk on disposal. A local-gamer-ID collision bug could
occur specifically after a remove-then-add sequence, because the ID was derived from the
live collection's current size rather than a monotonically increasing counter --- fixed with a
dedicated counter that never reuses a value. And application data sent before a peer's
handshake had fully completed was silently and permanently dropped, despite being reachable
through the ordinary public send API --- fixed with a small, bounded, per-session queue that
flushes once both sender and target resolve a real wire ID, then further hardened in a later
pass after a follow-up audit found the queue itself was never purged when a gamer left,
disconnected, or a host migration occurred. The bound is a real, fixed constant --- 64 pending
application-data sends per session, oldest evicted first once a 65th arrives before the same
peer's handshake completes --- with every eviction, whether from overflow or from the later
purge-on-disconnect fix, counted through an internal diagnostic counter
(\texttt{CNA::Internal::Net::ENetBackend::GetDroppedAppDataCount()}) rather than only logged;
this counter lives on the internal backend class, not on \cnaclass{NetworkSession}'s own public
surface, so it is a tool for this project's own test suite and diagnostics, not something
ordinary game code reads directly --- but it is worth knowing exists if you are ever debugging
a mysteriously missing early packet during connection setup.

A fifth, structurally distinct bug is worth naming on its own, since it lives one level below
this namespace rather than inside it: every \texttt{GamerCollection<T>} specialization ---
which includes \cnaclass{NetworkSession}'s own \texttt{AllGamers} / \texttt{LocalGamers} --- had
an enumerator whose \texttt{MoveNext()} unconditionally dereferenced the underlying collection
pointer after \texttt{Dispose()}, with no guard matching the one \texttt{getCurrent()} already
had for exactly that situation.

\begin{lstlisting}
// Worked example: the exact repro that surfaced the bug -- Dispose(),
// then MoveNext(), with no intervening getCurrent() call at all.
auto it = session->getAllGamersProperty().GetEnumerator();
it.Dispose();
it.MoveNext();   // used to dereference a null collection_ pointer here
\end{lstlisting}

Fixed by giving \texttt{MoveNext()} the identical \texttt{nullptr} check
\texttt{getCurrent()} already performed, throwing the same
\texttt{System::ArgumentOutOfRangeException} rather than crashing --- a small fix, but one that
existing tests had never actually exercised, since every prior test disposed an enumerator and
then called \texttt{getCurrent()}, never \texttt{MoveNext()}, on the disposed instance.

An earlier cross-module audit flagged a real leak against
\texttt{NetworkSession::BeginCreate}.  It is closed at the pin: every matching \texttt{End*}
deletes the action before clearing the shared slot, and a live-instance regression test covers a
complete create cycle.  The investigation still exposes a genuinely surprising architectural
constraint worth preserving here.

\subsection{One \texttt{Begin*} / \texttt{End*} slot, process-wide}

\texttt{activeAction\_}, the internal object every \texttt{BeginCreate} / \texttt{BeginFind} /
\texttt{BeginJoin} overload allocates to carry state across to its matching \texttt{End*} call,
is declared \texttt{static} on the class --- not a per-instance member:

\begin{lstlisting}
NetworkSession::NetworkSessionAction* NetworkSession::activeAction_ = nullptr;
\end{lstlisting}

This means exactly one asynchronous session-lifecycle operation may be in flight at a time,
\emph{across the entire process}, regardless of how many \cnaclass{NetworkSession} instances a
game happens to be juggling --- not one slot per session. Every \texttt{Begin*} overload
enforces this directly, throwing rather than silently queuing or overwriting a still-pending
action:

\begin{lstlisting}
if (activeAction_ != nullptr || activeSession_ != nullptr)
    throw /* ... */;
activeAction_ = new NetworkSessionAction( /* ... */ );
\end{lstlisting}

Under the ordinary, correctly-paired usage every synchronous \texttt{Create()} /
\texttt{Find()} / \texttt{Join()} convenience wrapper already follows --- call the matching
\texttt{End*} once the operation completes --- this is not a leak at all:
\texttt{EndCreate} / \texttt{EndFind} / \texttt{EndJoin} each call \texttt{delete activeAction\_;
activeAction\_ = nullptr;} exactly once, real and verified by reading each \texttt{End*}
directly. The single-slot design is a real, deliberate constraint worth knowing on its own
terms even apart from the leak question: a game that tries to \texttt{BeginFind} a second
session discovery while an earlier \texttt{BeginCreate} on a \emph{different}
\cnaclass{NetworkSession} object has not yet been matched with its own \texttt{EndCreate} will
hit this same guard and throw, a real cross-instance interaction a reader might not expect from
an API that otherwise treats each \cnaclass{NetworkSession} as independent.

The dated \texttt{NEXTaudio.md} note predates the Net fix.  Current source comments identify the
original defect precisely: each \texttt{End*} used to assign \texttt{nullptr} without first
deleting the heap-allocated \texttt{NetworkSessionAction}.  The pinned test
\texttt{CreateDoesNotLeakNetworkSessionAction} compares the live-instance count before and after
a full \texttt{Create()} cycle.  Separate LeakSanitizer reports caused by tests that call
\texttt{Dispose()} but never \texttt{delete} the caller-owned session pointer remain test-code
cleanup findings; they must not be used to reopen the corrected library leak.

\subsection{A worked example: the double-\texttt{Dispose()} use-after-free, traced in full}

The one-sentence summary above --- ``\texttt{Dispose()} could be called twice and use-after-free
as a result'' --- undersells a genuinely instructive root-cause chain, found by an
AddressSanitizer run that had nothing to do with networking at all (it was closing out an
unrelated content-pipeline fuzz-testing task; this was simply the first time the whole suite had
ever run under \texttt{-DCNA\_SANITIZE=address,undefined}). \texttt{ASan} reported a real
\texttt{heap-buffer-overflow} inside \texttt{NetworkSession::Dispose()}'s own
\texttt{ClearPacketQueue()} loop, surfaced by a real test
(\texttt{ENetBackendTest.DisposeDisconnectsConnectedPeersPromptlyInsteadOfWaitingForTimeout})
that was not written to test this bug at all --- it existed to confirm prompt disconnect
behavior, and triggered the crash purely as a side effect of its own fixture's cleanup:

\begin{lstlisting}
// From modules/net/tests/CNA/Internal/Net/ENetBackendTests.cpp:
{
    SystemLinkSessionFixture host;               // constructs a real NetworkSession
    host.session->Dispose();                     // (1) explicit Dispose(), by the test itself
}   // (2) fixture destructor unconditionally calls session->Dispose() AGAIN, no guard
\end{lstlisting}

Tracing exactly what the first \texttt{Dispose()} call does explains why the second one
crashes rather than being a harmless no-op. \texttt{Dispose()} clears each local gamer's packet
queue (fine --- the gamer object is still alive at this point), tears down the transport, then
calls \texttt{ownedGamers\_.clear()} --- which, since \cnaclass{NetworkSession} is the sole owner
of every gamer it created (the chapter's own ownership split, above), genuinely
\emph{destroys} the local \texttt{LocalNetworkGamer} object right there. Nothing at that point
prunes the now-dangling raw pointer out of \texttt{localGamers\_} / \texttt{allGamers\_} ---
only \texttt{RemoveGamer()} does that pruning, and \texttt{Dispose()} never calls it. The first
call finishes by setting \texttt{isDisposed\_ = true}, but by then the damage is already
structural: the collections still list a pointer to an object that no longer exists. The second
\texttt{Dispose()} call --- triggered by the fixture's own destructor, with no
\texttt{isDisposed\_} check of its own --- re-enters \texttt{ClearPacketQueue()}'s loop over
\texttt{localGamers\_}, dereferences that same dangling pointer, and \texttt{ASan} catches the
read.

This is a genuine API-contract violation, not merely a test-fixture quirk worth special-casing
away: real XNA's own \texttt{IDisposable} contract requires \texttt{Dispose()} itself to
tolerate being called more than once, and any real caller pairing an explicit
\texttt{Dispose()} call with a second, RAII-driven cleanup path --- exactly this fixture's own
shape, and a perfectly ordinary one --- would hit the identical crash outside of a test binary.
A follow-up, independent audit (\texttt{audit\_net.md}'s own Critical finding~1) went further
and found the bug was not even confined to the double-call case: after \emph{only one}
\texttt{Dispose()} call, a caller could already obtain a gamer collection containing the same
kind of dangling pointer, simply by reading \texttt{PreviousGamers} (which
\texttt{RemoveGamer()} populates with a departing --- possibly locally-owned --- gamer's raw
pointer rather than dropping it) or \texttt{Host} (a raw \cnaclass{NetworkGamer} pointer with the
identical exposure).

The fix closes both the double-call guard and the single-call exposure together, deliberately
placed inside \texttt{Dispose()} itself rather than pushed onto every caller to remember:

\begin{lstlisting}
void NetworkSession::Dispose()
{
    if (isDisposed_) return;   // the guard the destructor already had; Dispose() lacked it
    // ... existing teardown ...
    ownedGamers_.clear();
    // Defense-in-depth (audit_net.md Critical finding 1): a *single* Dispose() call
    // already left dangling pointers reachable through public collection properties.
    localGamers_.Clear();
    remoteGamers_.Clear();
    allGamers_.Clear();
    previousGamers_.Clear();   // RemoveGamer() populates this with raw pointers too
    host_ = nullptr;
    isDisposed_ = true;
}
\end{lstlisting}

Three new regression tests close the gap the original crash slipped through (every prior test
disposed an enumerator or session and then called a getter, never re-entered
\texttt{MoveNext()} / \texttt{Dispose()} on an already-disposed instance):
\texttt{DisposeCalledTwiceDirectlyIsSafeAndIdempotent} asserts \texttt{EXPECT\_NO\_THROW} on both
calls with every gamer count still reading zero afterward; the aptly-named
\texttt{DisposeClearsPreviousGamersSoNoDanglingPointerIsObservable} removes an owned local gamer
first (so it genuinely migrates into \texttt{PreviousGamers}), then disposes \emph{once}, and
confirms all four gamer-collection properties read back empty rather than exposing the dangling
pointer; and \texttt{DisposeClearsHostProperty} confirms \texttt{Host} reads back
\texttt{nullptr}. The original ASan repro test needed no changes at all --- it already exercised
exactly this path, and now serves permanently as the regression guard: a full
\texttt{-fsanitize=address} rerun of the exact same test reports zero heap-buffer-overflow or
use-after-free errors.

\section{A worked example: the \texttt{maxGamers} argument that this chapter's own opening example silently ignores}

The very first worked example in this chapter, at the top, passes \texttt{maxGamers = 8} to
\cnaclass{NetworkSession}'s own \texttt{Create()} --- and that value is never actually used.
This is worth correcting explicitly, in its own section, rather than quietly editing the
earlier example away, because the reason is a genuine, verified piece of API fidelity rather
than a bug in this port, and it is exactly the kind of trap a reader porting real game code
would otherwise hit directly. \texttt{Create(NetworkSessionType, int maxLocalGamers, int
maxGamers)}, the simplest three-argument overload, forwards straight into the matching simplest
\texttt{BeginCreate}
overload --- whose own \texttt{maxGamers} parameter is unused in its entire body, marked with an
explicit comment at the call site: FNA itself never uses this parameter in this overload,
preserved as-is. The reason is not merely that this one overload discards it: the internal
\texttt{NetworkSessionAction} object every \texttt{BeginCreate} overload constructs to carry
state across to the matching \texttt{EndCreate} call has no field to store a caller-supplied
\texttt{maxGamers} value in at all --- so \emph{every} \texttt{Create} / \texttt{BeginCreate}
overload in the whole class, not just the simplest one, ultimately reaches the same
\texttt{EndCreate}, which constructs the real \cnaclass{NetworkSession} with a literal, hardcoded
\texttt{69} in the \texttt{maxGamers} constructor slot, with its own comment recording exactly
why: \emph{``FNA hardcodes 69 here instead of forwarding the caller's original maxGamers
argument.''} \texttt{69} is not an arbitrary placeholder --- it is independently documented,
outside this project, as genuine historical Xbox~360/Windows XNA~4.0 behavior for this exact
call shape, which is why the prior audit that reconfirmed this quirk concluded it should be
preserved rather than corrected to ``do what the parameter name suggests.''

\begin{lstlisting}
// Worked example: the maxGamers argument is accepted, validated in some
// overloads, and then never actually reaches the constructed session --
// getMaxGamersProperty() below reads back 69, not 8, on every path.
NetworkSession* session = NetworkSession::Create(
    NetworkSessionType::SystemLink, /* maxLocalGamers */ 1, /* maxGamers */ 8);
assert(session->getMaxGamersProperty() == 69);  // not 8 -- real, preserved XNA behavior
\end{lstlisting}

A game genuinely relying on a specific room-size cap --- to size a lobby UI, or to refuse a
ninth join attempt at the application level rather than trusting the framework to enforce it ---
cannot use this parameter for that purpose on any \texttt{Create} / \texttt{BeginCreate} overload
in this class, and must track and enforce its own desired capacity separately, exactly as a game
targeting real Xbox~360-era XNA~4.0 would have had to.

\section{Two asymmetries, preserved rather than symmetrized}

Two small API details are worth knowing specifically because they look like bugs and are
not: \cnaclass{PacketReader} / \cnaclass{PacketWriter} read and write \cnaclass{Color} values
asymmetrically --- as four floats on read, four bytes on write --- a known, preserved mismatch
matching FNA's own upstream behavior exactly rather than a CNA-introduced inconsistency. And
\cnaclass{NetworkSessionProperties}, a C\texttt{++} translation of a C\# list type whose
single mutable indexer cannot distinguish a read access from a write access the way C\#'s
separate getter/setter accessors can, triggers FNA's own ``auto-append past the end'' growth
behavior on \emph{any} out-of-range access through its mutable overload, not only on an
explicit write --- a documented, accepted structural deviation; a \texttt{const} reference
gives genuinely bounds-checked reads instead.

\begin{lstlisting}
// Worked example: interoperating with Write(Color)'s actual four-byte
// representation. ReadColor() is NOT its inverse: it expects 16 bytes.
PacketWriter writer;
writer.Write(playerTintColor);      // 4 bytes on the wire
localGamer->SendData(writer, SendDataOptions::Reliable);

// On the receiving end, in the game loop:
std::vector<SharpRuntime::bytecs> rgba(4);
NetworkGamer* sender = nullptr;
while (localGamer->getIsDataAvailableProperty())
{
    const int count = localGamer->ReceiveData(rgba, sender);
    if (count == 4)
    {
        Color tint(rgba[0], rgba[1], rgba[2], rgba[3]);
    }
}
\end{lstlisting}

This example deliberately uses the byte-vector receive overload too.  The
\texttt{ReceiveData(PacketReader\&, NetworkGamer*\&)} overload does fill and rewind its
reader when a packet is available, but --- preserving an FNA bug --- its local length variable
is never updated and it returns zero even after consuming a packet.  Consequently a loop whose
condition is \texttt{ReceiveData(reader, sender) > 0} never executes its body.  If a port must
use that overload, gate each call with \texttt{IsDataAvailable} and treat the reader contents,
not the return value, as the successful-delivery signal; for new code, the vector overload above
has an honest byte count.

\section{\cnaclass{SendDataOptions}: five values, three real delivery guarantees}

The Color-asymmetry worked example above uses \texttt{SendDataOptions::Reliable} without
explaining what that value actually buys over the remaining four --- worth doing directly, since
\cnans{Microsoft::Xna::Framework::Net}'s own \cnaclass{SendDataOptions} enum has five values but
\texttt{ENetBackend} genuinely only has three distinct delivery behaviors underneath them, for a
real, documented reason rooted in what ENet itself can express. Reading
\texttt{NetPacketCodec::SendDataOptionsToEnetFlags()} directly, and the doc comment right above
it in \texttt{NetPacketCodec.hpp}, shows the exact mapping:

\begin{lstlisting}
uint32_t NetPacketCodec::SendDataOptionsToEnetFlags(SendDataOptions options)
{
    switch (options)
    {
        case SendDataOptions::None:
            return ENET_PACKET_FLAG_UNSEQUENCED;
        case SendDataOptions::InOrder:
            return 0;
        case SendDataOptions::Reliable:
        case SendDataOptions::ReliableInOrder:
        case SendDataOptions::Chat:
            return ENET_PACKET_FLAG_RELIABLE;
    }
    return ENET_PACKET_FLAG_RELIABLE;
}
\end{lstlisting}

\texttt{None} requests ENet's \texttt{UNSEQUENCED} flag explicitly --- true best-effort, packets
may be dropped or arrive out of order, the only \cnaclass{SendDataOptions} value that does.
\texttt{InOrder} maps to a bare \texttt{0}: ENet's plain unreliable send is already sequenced
per-channel and never delivered out of order on its own, so the empty flag set is an exact match
for ``in order, but not guaranteed to arrive'' rather than an oversight. The interesting case is
the other three: \texttt{Reliable}, \texttt{ReliableInOrder}, and \texttt{Chat} all collapse onto
the identical \texttt{ENET\_PACKET\_FLAG\_RELIABLE} flag, because ENet's reliable delivery
channel is \emph{inherently} ordered --- there is no ENet primitive for ``guaranteed to arrive,
but not necessarily in order,'' the specific combination real XNA's own \texttt{Reliable} (as
opposed to \texttt{ReliableInOrder}) documents. Rather than fake an unordered-reliable mode with
extra bookkeeping ENet does not need, the mapping accepts the collapse: any game that requests
plain \texttt{Reliable} gets \texttt{ReliableInOrder}'s exact behavior for free, a strictly
stronger guarantee than what it asked for, never a weaker one. \texttt{Chat} is a separate,
explicit judgment call rather than a parity match --- FNA itself never implements real delivery
for this value at all, so there is no existing behavior to preserve; routing it onto
\texttt{RELIABLE} instead of silently dropping chat text was a deliberate choice, confirmed by a
real, dedicated test:

\begin{lstlisting}
// From the real test suite (NetPacketCodecTests.cpp) -- confirms all five
// mappings directly, including the three-way collapse onto RELIABLE.
EXPECT_EQ(SendDataOptionsToEnetFlags(SendDataOptions::None),
          static_cast<uint32_t>(ENET_PACKET_FLAG_UNSEQUENCED));
EXPECT_EQ(SendDataOptionsToEnetFlags(SendDataOptions::InOrder), 0u);
EXPECT_EQ(SendDataOptionsToEnetFlags(SendDataOptions::Reliable),
          static_cast<uint32_t>(ENET_PACKET_FLAG_RELIABLE));
EXPECT_EQ(SendDataOptionsToEnetFlags(SendDataOptions::ReliableInOrder),
          static_cast<uint32_t>(ENET_PACKET_FLAG_RELIABLE));
EXPECT_EQ(SendDataOptionsToEnetFlags(SendDataOptions::Chat),
          static_cast<uint32_t>(ENET_PACKET_FLAG_RELIABLE));
\end{lstlisting}

The practical upshot for a game being ported: choosing between \texttt{Reliable} and
\texttt{ReliableInOrder} on this backend is purely a documentation choice about intent, not a
performance or behavior trade-off the way it would be on a transport that actually offers both
independently --- pick whichever name more accurately describes what the calling code needs, since
both compile down to identical wire behavior today. A future backend swap (a raw UDP
implementation, say, that genuinely can deliver reliable-but-unordered) is the only scenario where
the distinction would start to matter again.

\section{Status}

The namespace's own primary hardening plan is fully checked off. Three independent
post-completion audits each found that an earlier ``done'' declaration had overstated
reality --- mostly around Avatar-adjacent rendering concerns (Chapter~\ref{ch:avatar}) rather
than this namespace's own networking logic --- and the specific networking-side gap those
audits did surface (the pre-handshake application-data lifecycle above) is now fixed and
test-confirmed, with a fourth, final audit confirming the networking side as genuinely
settled. The current net module contains 293 statically counted test definitions. A project-wide
total would additionally need a named configuration and counting unit, as
Chapter~\ref{ch:counting-discipline} explains. Cross-machine delivery is proven by an
out-of-process two-instance ENet harness, while
discovery, codec hardening, policy, simulated latency/loss, and the deliberately preserved
\cnaclass{Color} asymmetry have focused tests. The history's more durable lesson is that each
``final'' audit found a different layer---lifetime, protocol, hostile input, or process
interaction---so a green feature checklist was repeatedly treated as a new claim to attack,
not as proof that further audit was unnecessary.
